\documentclass[11pt,a4paper]{article}
\usepackage[utf8]{inputenc}
\usepackage[T1]{fontenc}
\usepackage{amsmath,amssymb}
\usepackage{booktabs}
\usepackage[margin=2.4cm]{geometry}
\usepackage{hyperref}
\usepackage{xcolor}
\hypersetup{colorlinks=true,linkcolor=blue!50!black,urlcolor=blue!50!black}
\usepackage{titlesec}
\titleformat{\section}{\large\bfseries}{\thesection.}{0.6em}{}
\setlength{\parskip}{0.5em}
\setlength{\parindent}{0pt}

\title{\bfseries The Query Cannot See the Question\\[0.3em]
\large A short convolution's reach decides which part of a question\\
conditions retrieval from a model's memory}
\author{Maximiliano Speranza\\
\small Independent Researcher, Buenos Aires, Argentina\\
\small \href{https://orcid.org/0009-0005-0413-8554}{ORCID 0009-0005-0413-8554}}
\date{September 2026}

\begin{document}
\maketitle

\begin{center}
\fbox{\begin{minipage}{0.92\textwidth}
\small
\textbf{In one line.} In linear-attention and state-space models there is a part of every question
that the internal search \textbf{literally cannot see}, and what falls outside does not degrade
gracefully: \textbf{it disappears exactly}. The model then answers confidently from the part it can
still see.
\end{minipage}}
\end{center}

\section{What it is}

A model with memory has to do two things that are usually measured together and are in fact separate.
It has to \emph{find} the right entry, and it has to \emph{know} when the entry it is looking for is
not there. The second one is what separates a memory from a confabulator.

What we found is that the second one can fail for a purely mechanical reason, upstream of everything
the model learns. The query used to search the memory is not built from the whole question. It is
built by a \textbf{short causal convolution} over the hidden states, and that convolution has a fixed
reach of a few tokens.

If part of the question falls outside that reach, it does not influence the search a little.
\textbf{It does not influence it at all.}

\section{How it works}

The default kernel size of that convolution is \textbf{4}. That means the query at a given layer is a
function of the current token and \textbf{three predecessors}, and of nothing else.

That number is usually inherited rather than measured. Mamba-2 uses 4 by default. The reference
implementation of Qwen3-Next exposes \texttt{linear\_conv\_kernel\_dim}, documented as ``kernel size
of the convolution used in linear attention layers'', also defaulting to 4. A reach of three tokens is
not an exotic configuration, it is the factory setting in deployed models.

\subsection*{The concrete failure it produces}

Consider a question of the form \emph{``what is the \textbf{[relation]} of \textbf{[entity]}?''}. Read
from the end, the \textbf{entity} sits 1 token away and the \textbf{relation} 3. With a reach of 2,
the relation falls \textbf{exactly one token outside}, deterministically, in 100\% of queries.

The model therefore searches \textbf{using the entity alone}. And there the failure that most closely
resembles a real hallucination appears: asked about a relation that was \textbf{never stated}, for an
entity that \textbf{was}, it does not abstain. It retrieves the neighbouring fact, the one it does
have stored for that entity, and delivers it confidently.

It was not ignoring the relation. \textbf{It could not see it.}

\section{How it was found}

The intervention is simple enough to repeat on any architecture, and it requires no training:

\begin{center}
\fbox{\begin{minipage}{0.85\textwidth}
\centering\textbf{Change one token of the question and see whether the search moves.}
\end{minipage}}
\end{center}

Pick a read position, replace \textbf{a single token} at distance $d$ behind it, and measure how much
the convolution output moved at that position.

\subsection*{Measurement 1 · the real reach is not the nominal one}

Before measuring the effect one has to measure the window, because a weight can be zero, in which case
the effective reach is smaller than the nominal one. In \texttt{mamba-130m-hf}:

\begin{center}
\begin{tabular}{crrrr}
\toprule
layer & \multicolumn{4}{c}{$\max|\text{weight}|$ per tap, oldest to current} \\
\midrule
0  & \textbf{0.000000} & 0.933774 & 1.396596 & 3.513000 \\
1  & \textbf{0.000000} & 0.599272 & 0.928984 & 1.315802 \\
12 & \textbf{0.000000} & 0.309363 & 0.901561 & 0.981588 \\
23 & \textbf{0.000000} & 0.128891 & 0.611719 & 1.090700 \\
\bottomrule
\end{tabular}
\end{center}

\textbf{The oldest tap is exactly zero in all 24 layers.} Not small: zero, across all 1{,}536 entries
of each layer. The effective reach of this checkpoint is \textbf{2 tokens}, not 3.

We report the measurement and not a cause. A weight that is exactly zero across $24 \times 1536$
entries does not come from a gradient, so a conversion artefact is plausible, but we did not verify
that and we do not claim it. The practical consequence, if it generalises, is that \textbf{the nominal
kernel is an upper bound on the reach and should be measured rather than assumed}.

\subsection*{Measurement 2 · the cut is exact, and the layer sees what the query does not}

Here is the central result, and the part that surprises. With the same intervention we measure two
things: how much the \textbf{query} moves, and how much the \textbf{layer output} moves.

\begin{center}
\begin{tabular}{crr}
\toprule
distance & query & layer output \\
\midrule
$d=1$ & 1.100457 & 0.074570 \\
$d=2$ & 0.579417 & 0.057475 \\
$d=3$ & \textbf{0.000000} & 0.017182 \\
$d=4$ & \textbf{0.000000} & 0.024114 \\
$d=5$ & \textbf{0.000000} & 0.014883 \\
$d=6$ & \textbf{0.000000} & 0.012871 \\
$d=7$ & \textbf{0.000000} & 0.017005 \\
\bottomrule
\end{tabular}
\end{center}

Two things, and the second is what makes this non-obvious.

\textbf{There is no decay.} The query's movement goes from a large value to \texttt{0.000000} in one
step, exactly where the measured reach ends.

\textbf{And the model does see that token.} The layer output moves at \emph{every} distance, because
recurrence carries the information forward. That is:

\begin{center}
\emph{The state sees the whole sequence; the query that reads it does not.}
\end{center}

That is why the problem is invisible from the outside. The model \emph{has} the information. What it
cannot do is use it to decide what to look for.

\section{What it is good for}

\subsection*{A diagnosis that costs nothing}

Changing one token and checking whether retrieval moves requires no training, no gradients and no
labels. If it does not move, that token is \textbf{invisible} to the search, and any confidence the
model expresses about it is \textbf{unearned}.

\subsection*{A fix that costs 0.12\%}

Widening the kernel of the query-forming convolution from 3 to 5 raises the reach from 2 to 4, and the
missing part enters the window. Measured on a small model trained from scratch, three seeds per
condition:

\begin{center}
\begin{tabular}{lcc}
\toprule
 & kernel 3 & kernel 5 \\
\midrule
correct abstention in the hard case & $0.5850$ -- $0.7349$ & $\mathbf{0.9931}$ -- $\mathbf{1.0000}$ \\
overall accuracy & $0.91$ -- $0.93$ & $\mathbf{0.988}$ -- $\mathbf{0.993}$ \\
\bottomrule
\end{tabular}
\end{center}

\textbf{With no overlap} between conditions, against a trivial floor of $0.4065$. And it costs
\textbf{1{,}024 parameters out of 865{,}651}, that is \textbf{0.12\%} of the model: two more taps, by
128 channels, by 4 blocks.

It is not capacity. It is access.

\subsection*{And it is not a corner case}

Measured over \textbf{33{,}585 questions} from four public corpora (SQuAD, Natural Questions, TriviaQA
and HotpotQA), between \textbf{90\% and 100\%} of real questions have their discriminating parts
further apart than the measured reach. And it \textbf{gets worse with difficulty}: on multi-hop
questions it reaches $1.0000$.

The configuration our synthetic task was built to study is the ordinary case.

\section{What this does NOT claim}

The limit is part of the result, so it is worth stating explicitly.

\textbf{In a deep model the window attenuates rather than blocks.} The cut is exact at layer 0, but
from layer 1 recurrence restores the signal attenuated. We measured the rate: $1.077$ and $1.028$ per
token of distance in 24- and 48-layer models, with $r = 0.978$ in both. Three times the parameters and
twice the depth leave the rate unchanged, which is what makes the attenuation a property of the
architecture rather than of a particular checkpoint.

\textbf{The behavioural effect is transient when there are layers to spare.} Fine-tuning a 24-layer
model on the same task, the difference exists and is large at step 100, and \textbf{disappears
entirely by step 400}. This is a negative result, it was pre-registered, and we report it.

\begin{center}
\fbox{\begin{minipage}{0.88\textwidth}
\centering
\textbf{Where there are no layers left to pay with, the window sets a ceiling.\\
Where there are, it sets a toll.}
\end{minipage}}
\end{center}

The strong claim therefore holds for \textbf{memory consulted from an early layer}, and not as a
behavioural prediction for a deep model trained to convergence.

\section{Where it applies, by scale}

\begin{center}
\begin{tabular}{lll}
\toprule
 & depth & what happens \\
\midrule
small model & 4 blocks, read at block 0 & \textbf{permanent ceiling} \\
\texttt{mamba-130m} / \texttt{370m} & 24 and 48 layers & \textbf{toll}, the effect recovers \\
frontier models & tens of layers, hybrid & \textbf{not measured} \\
\bottomrule
\end{tabular}
\end{center}

On frontier models the honest statement is this: \textbf{the mechanism is present} --- kernel 4 is the
default in Mamba-2 and in the linear-attention layers of Qwen3-Next --- but \textbf{whether it
produces failures in a large model trained to convergence we did not measure and do not claim}. Two
reasons push against it: if 24 layers suffice to pay the toll, 60 are likely to suffice more; and
frontier models are hybrid, where a full-attention layer does not have this limitation because its
query sees the whole sequence.

Where it is reasonable to expect it to matter in a large system is when \textbf{memory is consulted
from an early layer}, as in retrieval conditioned on early-layer states. There are no layers upstream
to pay the toll with, and the ceiling returns.

\vspace{1em}
\begin{center}
\fbox{\begin{minipage}{0.9\textwidth}
\textbf{The lesson that generalises.} The query used to search a memory has to be formed where the
whole question has already been seen. And it can be checked cheaply on any architecture, without
training anything: change one part of the query and see whether retrieval moves. If it does not, that
part is invisible.
\end{minipage}}
\end{center}

\end{document}
